
\setcounter{page}{200}
\section{Trace for projective morphisms}

In this section we show the compatibility between trace morphisms for finite morphisms $\operatorname{Trf}$ and projective space trace $\operatorname{Trp}_f$, in the setting of the residue isomorphism (Proposition 8.2). This compatibility allows us to extend the trace formalism to all projectively embeddable morphisms. The result is provisional; our ultimate goal (Chapter VII) is a trace theory for arbitrary proper morphisms.

Consider the composition of functorial maps:
\[
G^\bullet \xrightarrow{\psi_{s,f}}
\Rfst s_* s^b f^* G^\bullet
\xrightarrow{\operatorname{Trf}_s}
\Rfst f^* G^\bullet
\xrightarrow{\operatorname{Trp}_f}
G^\bullet
\]
We prove this composite is the identity.

\begin{proposition}
Let $Y$ be a locally noetherian prescheme, $X=\mathbb{P}_Y^n$, $f\colon X\to Y$ the projection, and $s\colon Y\to X$ a section of $f$. Then for all $G^\bullet\in D_{\mathrm{qc}}^+(Y)$, the composition above is the identity morphism.
\end{proposition}
\begin{proof}
The maps used are from Lemma 8.1, Proposition 6.5, and Proposition 4.3 respectively.

1) Both $\psi_{s,f}$ and $\operatorname{Trp}_f$ may be computed via Cartan--Eilenberg resolutions of $f^*(G^\bullet)=f^*G^\bullet\otimes\omega_{X/Y}[n]$. We may use a common resolution, reducing the problem to the case of a single quasi-coherent sheaf $G$ on $Y$. Then $f^*G\otimes\omega_{X/Y}$ is $\mathcal{Ext}^n(s_*\sheaf{O}_X,\,\cdot)$-acyclic and $R^nf_*$-acyclic, so it suffices to verify the composition
\[
G \to f_*\mathcal{Ext}^n_X(s_*\sheaf{O}_X,f^*G\otimes\omega_{X/Y})
\to R^nf_*(f^*G\otimes\omega_{X/Y})
\to G
\]
is the identity.

\setcounter{page}{201}
2) The functors involved are right exact and commute with direct sums. The statement is local on $Y$, so we may assume $Y$ affine and reduce to $G=\sheaf{O}_Y$. We must show
\[
\sheaf{O}_Y \xrightarrow{\alpha}
f_*\mathcal{Ext}^n_X(s_*\sheaf{O}_X,\omega_{X/Y})
\xrightarrow{\beta}
R^nf_*(\omega_{X/Y})
\xrightarrow{\gamma}
\sheaf{O}_Y
\]
is the identity. This composition defines a canonical section $\delta(s)\in\Gamma(Y,\sheaf{O}_Y)$; we need $\delta(s)=1$.

3) The construction is stable under arbitrary base change, as all functors are compatible with base extension. The section $s$ may be obtained by base extension from the diagonal section $\Delta\colon\mathbb{P}^n\to\mathbb{P}^n\times\mathbb{P}^n$ over $\operatorname{Spec}\mathbb{Z}$. Thus $\delta(s)=(p_2s)^*\delta(\Delta)$, and we reduce to proving $\delta(\Delta)=1$. Since $\Gamma(\mathbb{P}^n,\sheaf{O}_{\mathbb{P}^n})=\mathbb{Z}$, $\delta(\Delta)$ is an integer; it suffices to test at a closed point.

\setcounter{page}{202}
4) We verify for the standard section $T_1=\cdots=T_n=0$ of $\mathbb{P}_Y^n$ over any $Y$ that $\delta(s)=1$. The proof uses explicit Čech and Koszul complex calculations:
- Use the standard affine cover of $\mathbb{P}_Y^n$ and relative differential sheaf $\omega_{X/Y}$.
- Use the Koszul complex associated to the local coordinates $T_1/T_0,\dots,T_n/T_0$.
- The map $\alpha$ is defined via the fundamental local isomorphism and canonical pairing on exterior powers.

\setcounter{page}{203}
The basis element $e_1\wedge\cdots\wedge e_n$ of the Koszul complex maps to the standard differential cocycle $\tau=dt_1\wedge\cdots\wedge dt_n$. The complex morphism from the Koszul complex to the Čech complex sends cochains to fractions in the standard affine cover.

\setcounter{page}{204}
Evaluating this cocycle yields the identity element in $\Gamma(Y,\sheaf{O}_Y)$, completing the proof.
\end{proposition}

\begin{corollary}
The isomorphism $\gamma$ of Theorem 3.4 is independent of the choice of homogeneous coordinates on projective space.
\end{corollary}
\begin{proof}
The isomorphisms $\psi_{s,f}$ and $\operatorname{Trf}_s$ are coordinate-independent, hence $\operatorname{Trp}_f$ is also canonical.
\end{corollary}

\begin{remark}
This construction extends to arbitrary locally free sheaves $E$ on $Y$: set $X=\mathbb{P}(E)$, define $\omega_{X/Y}$ as the relative dualizing sheaf, and construct the trace map $\operatorname{Trp}_f$ and duality theorem exactly as for projective space.
\end{remark}

\begin{proposition}
Let $u\colon X\to Y$ be a finite morphism of locally noetherian preschemes, $f\colon\mathbb{P}_Y^n\to Y$ the projective projection, and form the Cartesian fibre product diagram. Then the trace morphisms satisfy canonical compatibility in a commutative diagram of functors on $D_{\mathrm{qc}}^+(Y)$.
\end{proposition}
